\documentclass[../../main2.tex]{subfiles}

\begin{document}

	% ======================================================================
	% 骨架文件：小节标题与追问链为终版；正文由后续会话填充。
	% 扩增模式标注：本章 = 有压（强制执行）+ 定量编码。
	% 原书材料接入点：
	%   - 原书 part3 ch4（法律）→ §12.1 变长编码 + §12.2 定量意义拆分接入
	%   - §12.2 是全书两半的接合点：ch7 的 C_A 公式 = 法律定责公式
	%   - §12.4 嵌套关系必须用集合符号写出，且必须是推导出来的
	% ======================================================================

	\chapter{通道四：法律（博弈结果的编码与定量方法）}
	\label{ch:law}

	\begin{motivation}
		上一章遗留的问题：权力需要被约束——需要「即使权力方想作恶也无法」的机制。更普遍地：动态博弈产生了结果，双方对结果有争议，谁说了算？\\
		\textbf{核心动机}：论证法律的双重身份——它既是通道（变长编码的规则系统），更是\textbf{第7章贡献分解的现实应用}：$C_A = P(D \mid A) - P(D \mid \neg A)$ 就是法律定责的公式。这一章是全书两半的接合点。\\
		\textbf{扩增模式}：有压——强制执行消耗资源；编码一旦确立，本身抗变化。
	\end{motivation}

	\begin{logicmap}
	\begin{tikzpicture}[node distance=0.9cm, every node/.style={draw, rectangle, rounded corners, align=center}]
	\node (q) {博弈结果有争议，谁说了算？};
	\node (e) [below=of q] {法律是变长编码（§12.1）};
	\node (m) [below=of e] {法律是 Part III 的现实应用（§12.2）};
	\node (d) [below=of m] {法律驱动其余三通道（§12.3）};
	\node (n) [below=of d] {四通道嵌套关系（§12.4）};
	\draw[-{Stealth}] (q) -- (e);
	\draw[-{Stealth}] (e) -- (m);
	\draw[-{Stealth}] (m) -- (d);
	\draw[-{Stealth}] (d) -- (n);
	\end{tikzpicture}
	\end{logicmap}

	% ==================================================================
	\section{法律是变长编码}
	\label{sec:law-encoding}

	\textcolor{gray}{\textit{【骨架 · 追问链（终版）】沿用原书：歧义成本 > 编码成本，所以条文必须冗长。为什么要冗长？→ 因为要覆盖边界情形——常见情形短编码（「杀人偿命」），罕见情形长编码（几十页的例外条款）。}}

	\textcolor{gray}{（正文待写：变长编码的日常证据（民法典厚度）；「法律不是道德，是道德高频部分的压缩包」；与第9章语言压缩层同构。）}

	\begin{readingnote}
		\textcolor{gray}{（待写：你有没有想过，为什么「不得杀人」四个字人人知道，死刑复核却要几百页卷宗？）}
	\end{readingnote}

	\paragraph*{逻辑衔接}
	\textcolor{gray}{（待写）编码讲了法律的「形」。但法律的「魂」不在编码——在它被迫给出的那个数字。下一节是全书两半的接合点。}

	% ==================================================================
	\section{法律是 Part III 的现实应用}
	\label{sec:law-quantitative}

	\textcolor{gray}{\textit{【骨架 · 追问链（终版）】本节是全书两半的接合点：第7章的 $C_A = P(D \mid A) - P(D \mid \neg A)$ 就是法律定责的公式。法律不是「规则体系」，是「被强制要求给出确定数字的系统」。为什么要强制给数字？→ 因为纠纷必须了结——不能说「有点像他的错」，必须说「70\% 他的责任」。}}

	\textcolor{gray}{（正文待写：法学从「故意/过失」推进到责任比例；这个压力迫使法律发展出社科里最精细的因果分解工具——条件概率框架在此被反向验证；「一旦你接受责任 $\propto P(\text{损害} \mid \text{行为}) - P(\text{损害} \mid \neg \text{行为})$，你就得到了社会科学因果推断的统一语言」。）}

	\begin{questionbox}
		\textcolor{gray}{（待写：反驳位——「法律定责是拍脑袋，哪有真算概率？」回答：显式概率化确实罕见，但举证责任、优势证据标准的实质就是门槛化的概率语言——粗粒化的条件概率差。）}
	\end{questionbox}

	\begin{reader_anticipation}
		\item \textcolor{gray}{（待写）条件概率定责的严格形式？→ 第7章 §7.1--7.3 与附录 A.2/A.3。}
	\end{reader_anticipation}

	\paragraph*{逻辑衔接}
	\textcolor{gray}{（待写）法律的定量内核亮出来了。但作为通道，它对其他三条通道做什么？下一节「法律如何驱动其余三通道」。}

	% ==================================================================
	\section{法律如何驱动其余三通道}
	\label{sec:law-drives}

	\textcolor{gray}{\textit{【骨架 · 追问链（终版）】法律改变模因传播的合法性、经济活动的产权、政治权力的边界。法律是第四通道还是前三通道的约束？→ 两者都是——这必须诚实说清，不许一笔带过。}}

	\textcolor{gray}{（正文待写：三组例子（言论边界/产权界定/违宪审查）；「通道与约束」的双重身份：约束面朝存量，通道面朝增量——预测时两种用法都有效。）}

	\paragraph*{逻辑衔接}
	\textcolor{gray}{（待写）四条通道全部登场。它们到底是什么关系？下一节「四通道的嵌套关系」给出全书推导链的收口——这次是推导出来的，不是断言的。}

	% ==================================================================
	\section{四通道的嵌套关系}
	\label{sec:law-nesting}

	\textcolor{gray}{\textit{【骨架 · 追问链（终版）】文化 $\supset$ 经济 $\supset$ 政治（按约束嵌套：无压 $\to$ +稀缺 $\to$ +身份），法律为三者提供定量框架。嵌套意味着什么？→ 中间层不能独立存在：没有不传播模因的经济、没有不交换的政治。}}

	\textcolor{gray}{（正文待写：集合符号大图 + 每层「加入的约束」回放；介观统一论三命题的陈述（只描述，不证明，指向附录 A.5）——「经济是稀缺化的模因；政治是人情化的经济；法律是博弈结果的条件概率编码」。\lowfreq{三个命题的严格推导见附录/待发表论文《介观统一论》，此处只作直觉呈现。}）}

	\paragraph*{逻辑衔接}
	\textcolor{gray}{（待写）分解完成了：四条通道、每条的定义、判据、贡献算法、彼此的嵌套。Part V「预测与循环」回答最后的问题：怎么把贡献表拼回一条时间线。}

	% ==================================================================
	\section{本章小结与衔接}
	\label{sec:law-summary}

	\textcolor{gray}{（正文待写：收束 Part IV——统一视角一节：同一套「认同 + 资源 + 传播」动力学在不同约束下的同一张脸。）}

	\paragraph*{逻辑衔接}
	\textcolor{gray}{（待写）第13章「从贡献表到未来走向」：拿着 $C_{\text{模因}}, C_{\text{经济}}, C_{\text{政治}}, C_{\text{法律}}$ 四个数字，说出「接下来会怎样」。}

\end{document}
